\section{Selección y adaptación del método}

La oficina de gestión utilizará dirección híbrida basada en PMBOK y construcción iterativa. La adaptación sigue el documento PMI de 2021: elección del enfoque, ajuste organizacional, ajuste al proyecto y mejora continua. El marco fundamenta la adaptación; las reglas siguientes son decisiones de Synaptix. T-9 desarrolla la gestión del proyecto y T-10 el flujo de desarrollo, pruebas y liberación. \parencite[pp.~1--7]{sd6-pmi-adaptacion}

En este formulario, Proyecto identifica la función de gestión dirigida por el Jefe de Proyecto durante implementación y por el Gerente de Servicio durante Operación. El Jefe de Proyecto tendrá dedicación exclusiva en implementación y facultades para comprometer a Synaptix en ejecución. El relevo documentará responsabilidades y pendientes, sin eliminar su participación contractual posterior. Los responsables técnicos aportarán estado y evidencia de su ámbito; la nómina, suplencias y acreditaciones se documentarán en SD12/T-8 y las dedicaciones en SD7/T-15. \parencite[cap.~19, pp.~33--34]{sd6-bt}

La planificación y consolidación de resultados usarán quincenas civiles: 1--15 y 16--último día, coordinadas con T-10 y el Comité de Proyecto. El expediente de adaptación registrará decisión, condición del caso, práctica, responsable, producto, comprobación y documentos afectados. Su revisión mensual evaluará esperas, defectos, riesgos y carga del CLIENTE. Una mejora interna conservará capacidad y requisitos; si modifica compromisos contractuales, seguirá CC-6.


\section{Instancias de gobierno y comunicación}

La Tabla~\ref{tab:t9-comites} conserva las instancias obligatorias del artículo 71. Proyecto preparará agenda, expediente y decisiones requeridas; la secretaría registrará acuerdos, responsables y plazos. El Comité de Operación desde el mes 13 prepara la transición, sin adelantar la fase de servicio de los meses 21--56. \parencites[art.~71, pp.~37--38]{sd6-ba}[RT-19.09, p.~34]{sd6-bt}

\begingroup
\setlength{\tabcolsep}{3pt}
\begin{longtable}{
L{\dimexpr.20\textwidth-2\tabcolsep\relax}
L{\dimexpr.16\textwidth-2\tabcolsep\relax}
L{\dimexpr.32\textwidth-2\tabcolsep\relax}
L{\dimexpr.32\textwidth-2\tabcolsep\relax}}
\caption{Participación y decisiones de gobierno}
\label{tab:t9-comites}\\
\hline
\EncabezadoTabla{Instancia} &
\EncabezadoTabla{Cadencia} &
\EncabezadoTabla{Participantes} &
\EncabezadoTabla{Decisión y evidencia}\\
\hline
\endfirsthead

\hline
\EncabezadoTabla{Instancia} &
\EncabezadoTabla{Cadencia} &
\EncabezadoTabla{Participantes} &
\EncabezadoTabla{Decisión y evidencia}\\
\hline
\endhead

Ejecutivo &
Mensual &
Patrocinador del CLIENTE, gerencia de Synaptix y Administrador del Contrato. &
Estrategia, escalamiento, cambios de alcance y riesgos mayores; acta y autorización aplicable.\\
\hline

Proyecto &
Quincenal &
Contraparte Técnica y Jefe de Proyecto. &
Cronograma, entregables, desviaciones y decisiones operativas; acta y tablero.\\
\hline

Arquitectura &
Mensual &
Arquitecto de Solución, Seguridad y referentes técnicos del CLIENTE. &
Arquitectura, deuda y riesgos técnicos; acta y registro de decisión.\\
\hline

Operación &
Mensual desde mes 13 &
Líder de Operación, mesa de servicio y Contraparte Técnica. &
Niveles de servicio, incidentes, problemas y mejora; acta y acciones.\\
\hline

Seguimiento &
Semanal &
Equipos de ambas partes. &
Impedimentos y compromisos; registro de coordinación.\\
\hline
\end{longtable}
\FuenteTabla{Bases Administrativas, art.~71, pp.~37--38; Bases Técnicas Transversales, RT-19.09, p.~34.}
\endgroup

En el arranque se designarán titulares y suplentes, disponibilidad y alcance de autoridad. La suplencia será escrita y conservará las competencias exigidas; no transferirá automáticamente facultades contractuales. Proyecto comprobará representación, registrará ausencias y coordinará cobertura y escalamiento. La Contraparte designará referentes por proceso sin presumir un área TI del CLIENTE.

La secretaría registrará ID, fecha, asistentes, agenda, versiones examinadas, acuerdos o desacuerdos, autoridad, acciones, responsables, plazos y enlaces documentales. Distribuirá el borrador dentro del primer día hábil y publicará el acta dentro de los dos días hábiles siguientes a la sesión. Las discrepancias quedarán registradas sin aplazar la publicación ni atribuir aprobación por silencio. Toda corrección conservará la versión anterior, motivo y comunicación; las autorizaciones firmadas sólo podrán modificarse por sus autoridades competentes. Proyecto dará seguimiento y cerrará acciones con evidencia y confirmación del responsable de recepción.

El registro \texttt{INT-6} identificará interesado/rol, proceso, autoridad, impacto, disponibilidad, suplente, canal, responsable Synaptix y esfuerzo de participación. Funcional coordinará consultas con administración, torre, portería, escuela y varadero, agrupándolas según decisiones requeridas y ventanas seguras. El levantamiento y la ficha de requisitos se desarrollan en T-10. La participación distinguirá administración funcional, comités, pruebas y aceptación; las cuatro HH mensuales previstas en SD3 no se usarán como límite universal. \parencite[numerales 13.1--13.3, pp.~28--30]{sd6-c07}

El espacio colaborativo mantendrá documentación, entregables, actas, riesgos y cambios versionados, con permisos mínimos por perfil. Proyecto actualizará el tablero al cambiar un estado y al cierre diario, verificando acceso del CLIENTE en cualquier momento. La Tabla~\ref{tab:t9-comunicaciones} establece las comunicaciones principales; entrega o recepción no equivalen a aprobación. \parencite[RT-19.05/06/08, pp.~33--34]{sd6-bt}

\begingroup
\setlength{\tabcolsep}{3pt}
\begin{longtable}{
L{\dimexpr.17\textwidth-2\tabcolsep\relax}
L{\dimexpr.27\textwidth-2\tabcolsep\relax}
L{\dimexpr.23\textwidth-2\tabcolsep\relax}
L{\dimexpr.33\textwidth-2\tabcolsep\relax}}
\caption{Comunicaciones y constancias}
\label{tab:t9-comunicaciones}\\
\hline
\EncabezadoTabla{Producto} &
\EncabezadoTabla{Emisor y destinatario} &
\EncabezadoTabla{Plazo o disparador} &
\EncabezadoTabla{Canal, acción y constancia}\\
\hline
\endfirsthead

\hline
\EncabezadoTabla{Producto} &
\EncabezadoTabla{Emisor y destinatario} &
\EncabezadoTabla{Plazo o disparador} &
\EncabezadoTabla{Canal, acción y constancia}\\
\hline
\endhead

Agenda y acta &
Proyecto/secretaría a participantes. &
Agenda: dos días hábiles antes; acta: dos después. &
Espacio colaborativo y aviso; preparación/revisión y registro de distribución.\\
\hline

Informe mensual &
Proyecto a Contraparte y Administrador del Contrato. &
Corte mensual acordado. &
Espacio colaborativo y aviso; revisión de avance y constancia de entrega.\\
\hline

Desviación &
Proyecto a Contraparte y Administrador del Contrato. &
Cinco días hábiles desde detección superior al 10\,\%. &
Comunicación formal; revisión del plan y constancia de recepción.\\
\hline

Cambio aprobado &
Proyecto a responsables afectados. &
Antes de ejecutar. &
Comunicación formal y versión publicada; aplicación y registro de distribución.\\
\hline

Suspensión/reanudación &
Proyecto a equipos y Contraparte. &
Aviso de marejada y posterior habilitación. &
Canal operativo registrado; suspensión/reprogramación y constancia.\\
\hline

Mantenimiento &
Proyecto/SRE al CLIENTE. &
Al menos diez días hábiles antes. &
Comunicación formal; coordinación y constancia de recepción.\\
\hline
\end{longtable}
\FuenteTabla{Bases Administrativas, arts.~20, 71 y 72, y Bases Técnicas Transversales, cap.~19.}
\endgroup


\section{Planificación, capacidad y configuración documental}

Proyecto verificará por paquete de SD7 alcance, criterio de aceptación, etapa, predecesoras, responsables y ventana. Las estimaciones se construirán desde actividades, con contraste UCP para los núcleos; CPM identificará ruta crítica y holguras, y PERT considerará tres duraciones. La red no admitirá ciclos y la nivelación usará disponibilidad real por persona, incluyendo revisión, gobierno y soporte. La fecha inicial será un supuesto hasta formalizarla; feriados, campañas y fechas náuticas se incorporarán antes de aprobar la línea base.

La preparación, asistencia y seguimiento se imputarán una sola vez. Se conservan como propuestas a conciliar con SD7/T-15 los expedientes \texttt{1.1.4.m.q} de meses 1--20, con 16 HH por quincena: Proyecto 8, Funcional 4 y Calidad 4; y \texttt{1.12.3.1.m} de 21--56, con 16 HH mensuales: gestión del servicio 4, Calidad 4 y Operación 8. No representan toda la capacidad de gobierno: los demás participantes y el Jefe de Proyecto durante Operación enlazarán sus horas a las asignaciones correspondientes. Una persona con varias funciones conserva una sola capacidad; una insuficiencia exige cobertura y nivelación antes de comprometer trabajo, especialmente en 13--15 y 19--20. La mesa mantendrá relevo para asistir a sesiones.

Los cálculos se reproducirán mediante Python 3 y registros CSV UTF-8 de SD7, conservando programa/entorno, entradas y resultados. El repositorio Git del CLIENTE versionará registros y documentos, y el espacio autenticado publicará el tablero HTML y descargas PDF/CSV. Las actualizaciones quedarán auditadas y se respaldará su restauración; las aprobaciones serán documentos firmados, no movimientos del tablero.

El registro maestro enlazará requisito y versión de fuente, CA, etapa, entregable, paquete, responsable, prueba, evidencia, resultado y decisión. Proyecto administrará su versión y los responsables actualizarán sus ámbitos, referenciando \texttt{ITEM-6}/\texttt{LIB-6} de T-10 sin duplicarlos. Se conservarán los 576 IDs y 27 complementos del caso de T-12; modificar una declaración exige evaluar sus cláusulas y justificar la nueva versión.

Cada línea base identificará alcance, programación y presupuesto autorizado. Antes de aprobarla se resolverán diferencias documentales y se comprobarán dependencias, capacidad y ventanas. La aprobación registrará autoridad competente, fecha y versión; Proyecto comunicará la vigencia y conservará el historial. Un compromiso futuro no se presentará como prueba ejecutada. \parencites[RT-19.01/03/05/08, pp.~33--34]{sd6-bt}[art.~50.2, p.~29]{sd6-ba}


\section{Procedimiento de control de cambios}

El registro \texttt{CC-6} conservará ID, solicitante, fecha, motivo, clasificación justificada, requisitos/paquetes afectados, alternativas, impactos, decisión, responsables, versiones y verificación. Los importes permanecerán en el expediente económico autorizado. \parencites[art.~72, p.~38]{sd6-ba}[RT-19.03, p.~33]{sd6-bt}

\begin{enumerate}
\item Proyecto y Funcional registran y clasifican la observación.
\item Un defecto o precisión sin cambio contractual se corrige bajo la autoridad técnica aplicable, conservando requisitos, recepción y controles de T-10. Si altera compromisos, se reclasifica como modificación contractual.
\item Para una modificación, Arquitectura coordina los impactos en alcance, plazo, costo, riesgo y servicio con los responsables afectados; Proyecto verifica capacidad, ventanas e hitos.
\item Se comprueban el límite acumulado del 20\,\% del valor original del Contrato y la prohibición de modificar objeto, naturaleza, garantías mínimas o cronograma obligatorio del artículo 17. El Comité Ejecutivo aprueba y ambas partes formalizan por escrito antes de ejecutar.
\item Proyecto comunica y actualiza las versiones afectadas; Calidad verifica el resultado y la Contraparte recibe cuando corresponda. Rechazo, aplazamiento y cierre conservan su evidencia.
\end{enumerate}

Aprobar una modificación no equivale a aceptar el producto. Proyecto bloqueará la ejecución contractual mientras falten aprobación o acuerdo escrito. Una modificación ejecutada sin autorización será de cargo y riesgo de Synaptix y no dará derecho a pago adicional: se registrará, comunicará al Administrador del Contrato y contendrá de forma segura, sin alterar retrospectivamente la línea base. \parencite[art.~72, numerales 1--5, p.~38]{sd6-ba}


\section{Medición de avance y recuperación}

El procedimiento \texttt{AV-6} fijará por paquete presupuesto, calendario y regla de reconocimiento. Los productos previstos dentro de un período mensual aplicarán 0/100: valor sólo al terminar y verificarse. Los paquetes de varios períodos usarán hitos verificables ponderados por el presupuesto asociado, con pesos que sumen uno. Los servicios periódicos reconocerán resultados del período, sin anticipar meses futuros. Calidad comprobará evidencia técnica; aceptación y pago conservarán sus condiciones propias.

Con pesos $w_{ik}$, verificaciones aprobadas $a_{ik}\in\{0,1\}$ y presupuesto $B_i$:
\[
q_i=\sum_k w_{ik}a_{ik},
\qquad
EV=\sum_i B_iq_i.
\]
La regla 0/100 usa un único resultado de peso uno. $PV$ procede de la programación presupuestada y $AC$ de costos reales conciliados; $SPI=EV/PV$ y $CPI=EV/AC$ usan igual fecha de corte y base monetaria. Un denominador cero se declara no calculable. Ningún trabajo se reconoce dos veces; los importes se conservan fuera de la oferta técnica. \parencites[p.~1]{sd6-nasa}[RT-19.06/07, p.~34]{sd6-bt}[art.~50.2, p.~29]{sd6-ba}

El informe mensual incluirá cronograma, avance físico/financiero, entregables, desviaciones, riesgos, incidencias y compromisos siguientes. El Jefe de Proyecto lo consolidará en implementación y el Gerente de Servicio en Operación. Se contrastará SPI con fechas y ruta crítica, y se informarán por separado calidad, niveles de servicio y aceptación.

Para RT-19.10 se conservará la regla de SD6:
\[
D_h=100\,\frac{|F_h^p-F_h^0|}{F_h^0-I_h^0}.
\]
$F_h^p$ es el término proyectado; $F_h^0$, el término de línea base; e $I_h^0$, el inicio más temprano del bloque de actividades específicas de producción, integración, verificación y recepción que culminan en el hito. SD7 fijará ese bloque y su duración positiva, sin actividades ajenas ni cambios posteriores para reducir el porcentaje. Las diferencias se expresarán en días corridos; un hito instantáneo utilizará la duración del bloque.

El registro mostrará versión, fechas, cálculo, adelanto/atraso, holgura, primera detección y fecha límite de aviso. Una desviación superior al 10\,\% se comunicará dentro de cinco días hábiles desde esa detección, con causa, impacto, responsable, plazo y plan de recuperación. El riesgo de incumplir un hito se escalará desde su detección, aunque no alcance el umbral. La fórmula es una convención de Synaptix; las bases establecen umbral y plazo. \parencite[RT-19.10, p.~34]{sd6-bt}

La recuperación verificará alternativas y capacidad antes de paralelizar; una nueva proyección no sustituirá la línea base sin autorización. SD8/T-16 aplicará ISO 31000:2018 con causa, evento, impacto, evaluación inicial/residual, respuesta, responsable, plazo y disparador. Proyecto revisará riesgos en cada Comité de Proyecto y ante incidentes o cambios, enlazando acciones y consumo de reserva una sola vez. Un riesgo materializado conservará su relación con la incidencia; no se aceptarán excepciones a requisitos obligatorios como riesgo residual. \parencites[RT-19.04, p.~33]{sd6-bt}{sd6-iso31000}


\section{Adquisiciones y ventanas operacionales}

El registro \texttt{ADQ-6} vinculará cada dependencia con SD7/T-11: responsable de provisión, especificación/cantidad, variante, interfaces, fecha necesaria, plazo de suministro, recepción y disponibilidad. Separará obras y hardware asignados al CLIENTE por el Caso 07 de bienes y servicios ofertados por Synaptix. Synaptix especificará y coordinará integración; los supuestos adicionales no eximirán obligaciones. Ante atraso se registrarán impacto, alternativa, responsable y paquete afectado antes de comprometer instalación o prueba. \parencite[cap.~11, pp.~26--27]{sd6-c07}

Se compararán compatibilidad, soporte, licencia, seguridad, entrega y sustitución. Arquitectura comprobará interfaces; Seguridad, licencias y tratamiento de datos; Calidad, recepción técnica. Las características se sustentarán en la variante real, con margen de suministro y recepción; los valores permanecerán en el expediente económico.

La subcontratación exigirá autorización escrita, control del límite acumulado del 40\,\% y restricciones de operación/custodia del núcleo crítico. Los terceros mantendrán los estándares exigidos y Synaptix sus responsabilidades; quienes accedan a datos serán declarados, auditados y sujetos al acuerdo aplicable. SD12 conservará proveedores y respaldos, sin presentar autorizaciones o cartas pendientes como obtenidas. \parencite[art.~73, p.~38]{sd6-ba}

Ante marejadas se suspenderá inmediatamente toda actividad programada, incluida la aislada o remota. Proyecto registrará y comunicará afectados, respuesta \texttt{PL-R01/R-01} y consumo de reserva; sólo reprogramará tras cesar la condición y comprobar autorización, capacidad y ventana segura. Se respetarán congelamiento total del 15 de diciembre al 15 de marzo, campañas de abril--mayo y septiembre--noviembre en lo que afecte al varadero, y congelamientos locales de las 14 fechas náuticas.

El mantenimiento programado se avisará con al menos diez días hábiles, conservando constancia de recepción; SRE verificará esa condición mediante T-10. Las reservas de SD7 no se duplicarán ni autorizarán trasladar hitos. Los incidentes reales seguirán su procedimiento sin habilitar trabajo programado durante restricciones. \parencites[cap.~10, restricción 12, p.~26; numeral 13.2, p.~29; cap.~15, RT-10.05, p.~34]{sd6-c07}[art.~20, p.~15]{sd6-ba}[RT-10.05, p.~22]{sd6-bt}


\section{Comprobación del método y cierre}

Antes de aplicar el método, Proyecto y Calidad comprobarán ambas rutas de CC-6, aprobación de una línea base, trazabilidad requisito--recepción, reproducción de un corte AV-6 y cálculo de desviación, y correspondencia entre carga y capacidad. Registrarán responsables, versiones, resultados y observaciones resueltas. Los ejemplos o simulaciones se identificarán como tales; no sustituirán ensayos técnicos de T-10 ni aceptación contractual.

El cierre por etapa conservará productos, observaciones resueltas y aceptación T-17/T-18. La conciliación de marcha blanca será diaria y el cambio oficial de autoridad esperará su acta. El expediente reunirá entregables, manuales, decisiones, riesgos residuales y pendientes con responsable y plazo. El relevo al Gerente de Servicio documentará recepción de registros y responsabilidades; las lecciones alimentarán la mejora y SD7 conservará calendario y responsables de reversibilidad y salida. \parencites[art.~17.3, pp.~12--13]{sd6-ba}[cap.~20, pp.~34--35]{sd6-bt}